\documentclass[12pt, a4paper, openany]{report}

\usepackage[left=3cm,top=3cm, bottom=3cm, right=4cm]{geometry}

\usepackage{secdot} % Dots in Section Numbers
\usepackage[utf8]{inputenc}
\usepackage[ngerman]{babel}
\usepackage{graphicx}
\graphicspath{ {images/} }

\usepackage{fancyhdr}
\usepackage{todonotes}
\usepackage{hyperref}

\usepackage{titlesec}
\titleformat{\chapter}[hang]{\Huge\bfseries}{}{0pt}{\Huge\bfseries}
\newcommand\frontmatter{ \cleardoublepage \pagenumbering{roman}}
\newcommand\mainmatter{ \cleardoublepage \pagenumbering{arabic}}
\newcommand\backmatter{ \if@openright \cleardoublepage \else \clearpage \fi }

\usepackage[autostyle]{csquotes}
\usepackage[
    backend=biber,
    style=alphabetic,
    sortlocale=de_DE,
    natbib=true,
    url=false,
    doi=true,
    eprint=false
    ]{biblatex}
\addbibresource{literatur.bib}

\pagestyle{fancy}
\fancyhf{}
\lhead{Jan van Dick}
\chead{Bachelorarbeit: \glqq Zweite Natur und Befreiung\grqq}
\rhead{\thepage}

\title{
    {Zweite Natur und Befreiung - zwischen Affirmation und Kritik}\\
    {\large Goethe Universität Frankfurt am Main}\\
    {\includegraphics{logo.png}}
}
\author{Jan van Dick}
\date{\today}

\begin{document}

\frontmatter
\maketitle

\chapter*{Abstrakt}
Friedrich Nietzsches These vom Tode Gottes, dass \glqq wir\grqq{} ihn getötet haben, ist mehr als bloße Negativität. 
Müsste der Mensch nicht selbst Gott werden, um ihn getötet zu haben?.\footfullcite[Vgl.][481]{nietzsche_morgenrote_1999} 
Und der Mensch ist Gott dadurch geworden, dass er ihn \textit{geschaffen} hat, dadurch nur konnte er ihn töten. 
In dem Tod Gottes, liegt, dass Gott selbst vom Menschen gesetzt, dass er Schein ist, der sich gegen ihn verselbstständigte.
Gott ist somit geistiges Produkt des Menschen, in dem der Geist selbst wieder zur Natur sich verkehrte.
Der Tod Gottes ist die Befreiung daraus.
Der \glqq Tolle Mensch\grqq{} erklärt aber zugleich, dass nur die wenigsten von dieser Tat wissen. 
Für die Einen ist der Tod Gottes, eine untergegangene Sonne, die Befreiung also wieder eine In-Natur-Verkehrtheit.
Nur für uns \glqq geborene Räthselrather\grqq, die den Tod Gottes im vollen Umfang begreifen, ist er ein \glqq neues offenes Meer\grqq\footnote{\citeauthor[][573]{nietzsche_morgenrote_1999}}.
Ist nun der Tod Gottes, die Befreiung aus der von ihm gesetzten Natur, das unbekannte, neue, offene, noch nie so offen gewesene Meer, oder ist er der Beginn des Wieder-in-Natur-Verkehrt-Seins, wie Christoph Menke es in der Analyse Hegels Begriffs der zweiten Natur beschreibt?\footfullcite[Vgl.][144]{menke_autonomie_2018}\\
Befreiung steht also zwischen Kritik (wieder-in-Natur-Verkehrtheit) und Affirmation (Einheit von Setzen und Sein).
Während Menke (und Hegel) Befreiung aus der zweiten Natur, nicht ohne eine neu hervorgebrachte zweite Natur, in welcher der Geist wieder in Natur verfällt lesen, erörterte ich die Frage, ob es in dem Motiv des \glqq neuen offenen Meeres\grqq{} und der \glqq gebornenen Rätselrather\grqq{} bei Nietzsche, eine Befreiung aus der fortwährenden In-Natur-Verkehrtheit geben kann.
Es ist die Frage nach einer dritten Befreiung neben der Befreiung aus der 1. und der 2. Natur. 
Die Antwort dazu wird sich in der Arbeit als Ergebnis des Unterschiedes der \glqq Erkennenden\grqq{} bei Nietzsche geben:
Zwar kann Befreiung aus der Natur nicht ohne Setzen einer zweiten Natur geschehen, aber die Befreiung, die sich in dem Bewusstsein ihrer eigenen Kritik vollzieht, ist zugleich über ihrer eigene Befreiung hinaus.
Die dritte Befreiung ist damit allerdings keine Befreiung aus der Befreiung.
Befrieung bleibt notwendig: der Mensch kann Freiheit nicht \textit{haben}, er muss sie immer wieder selbst hervorbringen.


\tableofcontents

\mainmatter

\chapter{Einleitung}
So, wie Friedrich Nietzsche in dem Aphorismus vom \glqq Tollen Menschen\grqq in \textit{Die fröhliche Wissenschaft} die These vom Tode Gottes entfaltet enthält sie weit mehr als die bloße Negation Gottes, hier steckt viel mehr ein dialektischer Begriff von (zweiter) Natur und Befreiung.
Diese beiden Begriffe möchte ich aus der sprachgewaltigen und metaphergeladenen Sicht Nietzsches, sowie der von Christoph Menke rekonstruierten Perspektive Hegels untersuchen.
Beide Begriffe entfalten sich in ihrem Doppelcharachter: zweite Natur - Kritik und Affirmation, Befreifung - Macht und Ohnmacht.
In der zweiten Natur schlägt Setzen in Sein um; hierin liegt zum Einen die Verwirklichung des Geistes (Affirmation), zum Anderen, der In-Natur-Verkehrtheit, der Tod des Geistes (Kritik).\footfullcite[Vgl.][145]{menke_autonomie_2018}
Ebenso verhällt es sich mit dem Begriff der Befreiung: 
Sie ist 1. \textit{Macht} des Geistes neues Hervorzubringen und die Notwendigkeit des Bestehenden zu durchbrechen, 
2. \textit{Ohnmacht} des Geistes, da die Befreiung sie nie abgeschlossen ist; die Macht der Befreiung ist die Schaffung der einer neuen zweiten Natur, aus der der Geist sich befreite.\footnote{\citeauthor[Vgl.][80]{menke_autonomie_2018}}\\
Eine ähnliche Struktur ist so auch bei Nietzsche zu finden.
Doch während bei Hegel Befreiung aus der zweiten Natur auf Grund der Endlichkeit des menschlichen Geistes wieder in Natur verfällt, scheint bei Nietzsche in der Metapher des neuen, offenen Meeres, die Perspektibe die ewige Wiederholung der In-Natur-Verkehrtheit zu überwinden.
Zugleich betonen sowohl Hegel, als auch Nietzsche Freiheit als nicht-gegeben: 
Freiheit wird bei beiden so gedacht, dass man sie \glqq nicht nur hat, sondern auch beständig  noch erwirbt und erwerben muss\grqq\footfullcite[][637]{nietzsche_morgenrote_1999}. Freiheit ist die Befreiung aus der jeweiligen Unfreiheit.\footfullcite[Vgl.][227]{adorno_negative_dialektik_2003}
Die Dialektik zwischen Freiheit und Notwendigkeit, Geist und Mechanismus, Endlichkeit und Unendlichkeit des Geistes, ist demnach Grundlage meiner Arbeit. Die Leitfrage der Arbeit also: \textit{gibt es in Nietzsches Philosophoe eine Möglichkeit der Überwenigdurch der Notwendigkeit-in-Natur-Verkehrtheit des Geistes?}\\
Aus der Position des oben erähnten \glqq Tollen Menschen\grqq{} liegt es nahe, diese Möglichkeit in dem Vollzug der Befreiung denken, in einer Form des Bewusstseins im Volzug der Handlung; 
eine Befreiung, die sich in dem Bewusstsein ihrer eigenen Kritik vollzieht.
Zu Fragen ist, ob dieses Bewussts-Werden eine Arbeit des einzelnen Individuums ist, oder ob es aus der Struktur des Allgemeinen, aus der zweiten Natur selbst, hervorgehen kann.





\chapter{Hauptteil}
\section{Zweite Natur bei Hegel}
\subsection{Die Sittlichkeit}
\subsection{Dialektik von Geist und Mechanismus}
\subsection{Dialektik von Kritik und Affirmation}
\section{Zweite Natur bei Nietzsche}
\subsection{Was bedeutet Schein?}
\subsection{Der Schauspieler und die Rolle}
\subsection{Gott ist tot!}
\section{Was heißt Befreiung aus der zweiten Natur}
\section{Das offene Meer?}

\chapter{Fazit}

\backmatter

\printbibliography
 
\end{document}
